% !TEX program = xelatex
% ============================================================
% 第 2 章《基本物理》§2-1、§2-3 内容层改造样张 —— 正文
%
% 约定：
%   * 原文一字未删：只插入分段、列表环境、\runin / \term / \keyidea /
%     \pullquote / \aside、导读、小结与示意图；
%   * 本文件中“新增文字”如下（零删除校验时整块剔除）：
%       \chapterintro、\readingpath、\runin、\enterm、\figsource，
%       summarybox / conceptmap / figure 环境，\label；
%   * 原文文字只换了呈现方式（\footnote -> \aside、段落切分），
%     校验脚本：python3 verify_nodel.py
% ============================================================

\chapter{基本物理}

\chapterintro{%
\textbf{一句话定位}：本章不推公式，只画一张“物理学的世界观地图”——物理学家在做什么、已经知道什么、还不知道什么。\par
\textbf{本章要点}：
\begin{itemize}
  \item 基本物理的目标，是把无穷多的现象归结为少量基本事物和相互作用。
  \item “理解”就像看懂一盘象棋：先学会规则，再谈走法；规则对不对，有三种检验办法。
  \item 物理学史是一连串“综合”：热与力学、电与磁与光、化学与原子行为。
  \item 1920 年后的两大变化：空－时取代绝对空间与时间，量子力学改写了小尺度上的规则。
  \item 量子力学的三大观念：测不准原理、不可精确预言、波粒二象性；量子电动力学把光与物质统一起来。
\end{itemize}
\textbf{与前后章的关系}：上承第 1 章“原子假说”的结论，下接第 3 章“物理学与其他科学的关系”；本章界定的“基本物理”是全书的纲领。\par
\textbf{本样张范围}：只改造 §2-1 与 §2-3 两节（§2-2 未收录），用来演示“导读—分段—列表—路标—图示—小结”的完整改造链条。
}

\readingpath{速读路径：§2-1 全节（本章纲领）→ §2-3 只看五个段内小标题与文末小结 → 回读加粗术语。精读路径：按正文顺序读，遇到“选读”标记的段落可先跳过，不影响主线；每节末尾的“小结”可当复查清单。}

% ============================================================
% §2-1 引言
% ============================================================
\section{引言}

\runin{本章要问什么} 在本章中，我们将考察有关物理学的最基本概念——即我们在目前所知道的事物的本性。这里将不去涉及我们如何知道所有这些观念是正确的那个认识过程，你们在适当的时候会学习到这些具体的细节。

\runin{从纷繁现象到基本事物} 我们在科学上所关心的事物，具有无数形式和许多属性。举例来说，假如我们站在岸边眺望大海，将会看到：这里有海水、拍击的浪花、飞溅的泡沫以及汹涌的波浪，还有太阳、光线、蔚蓝的天空、白云以及空气的流动——风；在海边有沙粒，不同色纹和硬度的岩石；在海里浮游着生物，此生彼灭；最后，还有我们这些站在海岸边的观察者；甚至还有幸福和怀念。

在自然界的其他场合，难道不也同样出现如此纷繁复杂的事物和影响吗？无论在哪里，到处都是这样错综复杂和变化无穷。好奇心驱使我们提出问题，把事物联系起来，而将它们的种种表现理解为：或许是由较少量的基本事物和相互作用以无穷多的方式组合后所产生的结果。

例如，

\begin{itemize}
  \item 沙粒和岩石是两回事吗？就是说，沙粒只不过是大量的细小石块吗？
  \item 月亮是不是一块巨大的岩石呢？如果我们了解岩石，是否就能了解沙粒和月亮呢？
  \item 风是否与海洋中的水流相似，就是一种空气的流动？
  \item 不同的运动有什么共同特征？
  \item 不同的声音有什么相似之处？
  \item 究竟有多少种颜色？
\end{itemize}

等等，等等。我们就是试图这样地逐步分析所有的事情，把那些乍看起来似乎不相同的东西联系起来，希望有可能\uwave{减少不同类}事物的数目，从而能更好的理解它们。

几个世纪以前，人们想出了一种部分解答这类问题的方法，那就是：\term{观察}，\term{推理}和\term{实验}；这些内容构成了通常所说的\term{科学方法}。在这里，我们将只限于对那些有时称之为\term{基本物理}中的基本观点，或者由于应用科学方法而形成的基本概念作一描述。

\runin{“理解”就像看懂一盘棋} 现在我们要问：所谓“理解”某种事情指的是什么意思？可以作一想象：组成这个“世界”的运动物体的复杂排列似乎有点像是天神们所下的一盘伟大的象棋\aside{（这里指的是国际象棋——译者注）}，我们则是这盘棋的观众。我们不知道弈棋的规则，所有能做的事情就是\uwave{观看}这场棋赛。当然，假如我们观看了足够长的时间，总归能看出几条规则来，这些弈棋规则就是我们所说的基本物理。但是，即使我们知道了每条规则，仍然有可能不理解为什么下棋时要走某一步棋，这仅仅是因为情况太复杂了，而我们的智力却是有限的。

如果你们会下棋，就一定知道，学会所有的规则是容易的，但是，要选择最好的一着棋，或者要弄懂别人为什么走这一着棋，往往就很困难了。在自然界里，也正是如此，而且只有更难一些。但是，至少我们能发现所有的规则。

实际上我们今天还没找到一切规则（时而会出现一些像弈棋中“以车护王”那样的情况，使我们仍然感到无法理解）。除此之外，我们确实能用已知规则来解释的事情也是非常有限的，因为几乎所有的情况都是极其复杂的，我们不能领会这盘棋中应用这些规则的走法，更无法预言下一步将要怎样。所以，我们必须使自己只限于弈棋规则这个比较基本的问题。

\pullquote{如果我们知道了规则，就认为“理解”了世界。}

\runin{检验规则的三种办法} 如果我们不能很好的分析这盘象棋游戏，那么又怎样来辨别我们“猜测”出的规则实际上是否正确呢？大致地讲，可以有三种办法。

\begin{enumerate}
  \item[第一，]可能有这种情况：大自然安排的，或者说我们将大自然安排的十分简单，只有少数几个组成部分，从而使我们能够正确地预测将要发生的事。在这种情况下，就能检验我们的规则是怎样起作用的。（在棋盘角落里可能只有少数几个棋子在移动，所以我们能够正确地解决。）

  \item[第二种]检验规则的好办法是，利用那些由已知规则推导出来的一些较一般性的法则来检验已知规则本身。比如，象在棋盘中移动的规则是只许走对角线，因而我们可以推断，无论象走了多少步，它总是出现在红方块里。这样，即使不能领会细节，我们也总能检验有关象的走法的概念，只要弄清楚它是否一直在红方块里。当然，在相当长的时间里，它都将如此，直到突然发现它出现在黑方块里。

（显然，这时发生的情况是这个象被俘获了，另一个卒走过来成为皇后，红方块里的象就变成黑方块里的象。）

这也就是物理学中出现的情况，即使我们不能领会其中的细节，但是在相当长的时期内我们仍有在各方面都很好地起作用的规则；但是在某个时候，我们又会发现新的规则。\keyidea{从基本物理的观点来看，最有趣的现象当然是在那些新的场合——那些已知规则行不通的场合中所出现的现象，而不是在原有规则行得通的地方发生的现象！}这是我们发现新规则的一条途径。

  \item[第三个]鉴别我们的观念是否正确的方法比较粗糙，但或许是所有方法中最为有效的。这就是用粗略的近似方法来加以辨别。我们可能说不出为什么阿莱克因（Alekhine）\aside{世界著名弈棋手，系国际象棋大师。曾多次获得国际象棋世界冠军。——译者注}要走这步棋，但是我们或许能大致认为它或多或少地在调集一些棋子到王的周围来保护它。因为这是在这种情况下明摆着的事。同样，根据我们对这盘棋的理解，即使不能看出每一步棋的作用，也常常能对自然界多少有所理解。

\end{enumerate}


\begin{conceptmap}[示意：检验规则的三种办法]
\centering
\begin{tikzpicture}[font=\small,>=stealth,
  piece/.style={circle,draw,inner sep=1pt,font=\scriptsize},
  panelcap/.style={align=center,text width=3.6cm,font=\footnotesize,text=black!80}]
% ① 棋盘角落：棋子少，能精确预测
\begin{scope}
  \draw[step=.55,gray!55,thin] (0,0) grid (2.2,2.2);
  \node[piece,fill=black!75,text=white] at (.275,.275) {兵};
  \node[piece,fill=black!75,text=white] at (1.925,1.925) {王};
  \node[piece,fill=white] at (.275,1.925) {车};
  \node[panelcap,below=3mm] at (1.1,0) {(1) 角落里棋子少\\可以精确预测};
\end{scope}
% ② 象只走对角线：永远落在同色格
\begin{scope}[shift={(4.4,0)}]
  \foreach \x in {0,...,3} \foreach \y in {0,...,3} {
    \pgfmathparse{mod(\x+\y,2)==0 ? 1 : 0}
    \ifnum\pgfmathresult=1
      \fill[red!18] (\x*.55,\y*.55) rectangle ++(.55,.55);
    \fi
  }
  \draw[step=.55,gray!55,thin] (0,0) grid (2.2,2.2);
  \draw[-{Stealth[length=2mm]},thick,red!65!black,dashed]
      (.275,.275) -- (1.375,1.375) -- (1.925,1.925);
  \node[piece,fill=black!75,text=white] at (.275,.275) {象};
  \node[piece,fill=white] at (1.925,1.925) {象};
  \node[panelcap,below=3mm] at (1.1,0) {(2) 只走对角线\\永远落在同色格};
\end{scope}
% ③ 以车护王：粗略近似，把握意图
\begin{scope}[shift={(8.8,0)}]
  \draw[step=.55,gray!55,thin] (0,0) grid (2.2,2.2);
  \node[piece,fill=black!75,text=white] (k) at (1.1,1.1) {王};
  \node[piece,fill=white] (r1) at (.275,.275) {车};
  \node[piece,fill=white] (r2) at (1.925,1.925) {车};
  \draw[-{Stealth[length=2mm]},thick,red!65!black] (r1) to[out=40,in=220] (k);
  \draw[-{Stealth[length=2mm]},thick,red!65!black] (r2) to[out=220,in=40] (k);
  \node[panelcap,below=3mm] at (1.1,0) {(3) 以车护王\\粗略近似也管用};
\end{scope}
\end{tikzpicture}

\figsource{按 §2-1“三种检验办法”的文字描述绘制，非原书插图}
\end{conceptmap}


\runin{分类、综合与拼图} 人们首先把自然界中的现象大致分为几类，如热、电、力学、磁、物性、化学、光或光学、X射线、核物理、引力、介子等等现象。然而，这样做的目的，是将\uwave{整个自然界}看作是一系列现象的许多不同侧面。\keyidea{这就是今天基础理论物理面临的问题：发现隐匿在实验后的定律；把各类现象综合起来。}在历史上，人们总能做到这一点，但随着时间的推移，新的事实发现了；我们曾经将现象综合得很好，突然，发现了 X 射线，随后我们又融合了更多事实，但是又发现了介子。因此，在弈棋的任何一个阶段，看起来总是相当凌乱。大量事实被归并了，但总还有许多线索向一切方向延伸出去。这就是今天的状况，也就是我们将试图去描绘的现状。

历史上出现过的若干进行综合的情况有如下几个。

\begin{itemize}
  \item \textbf{首先，}是\uwave{热}与\uwave{力学}的综合，当原子运动时，运动得越是剧烈，系统所包含的热量就越多，这样，\uwave{热和所有的温度效应}可以用力学定律来说明。
  \item \textbf{另一个巨大的综合，}是发现了电、磁、光之间的联系，从而知道它们是同一件事物的不同方面，即今天我们称为\term{电磁场}的那个东西的不同表现。
  \item \textbf{还有一个综合，}是把化学现象、各种物质的各种性质以及原子的行为统一起来，这就是\term{量子化学}的内容。
\end{itemize}

\runin{世界由什么构成} 显然，现在的问题是：能不能继续把所有的事情都综合起来，并且仅仅发现这整个世界体现了\uwave{一件}事情的种种不同方面？无人知道答案如何，我们所知道的只是：这样做下去时，我们发现可以综合一些事实，随后又发觉出现了一些不能综合的事实。我们继续尝试这种拼图游戏。至于是否只有有限数量的棋子，甚至这场拼图游戏是否有底，当然不知道。除非有那么一天终于把拼图拼成了，否则我们就永远不会知道事情的究竟。在这里我们要做的是，看看那种综合已经到了什么程度，在借助于最少的一组原理来理解基本现象方面，现状又是如何。\keyidea{简言之，事物是用什么构成的？总共存在多少基本元素？}


\begin{summarybox}[§2-1 小结]
\begin{itemize}
  \item 基本物理要做的事：把无穷多的现象归结为少量基本事物与相互作用。
  \item “理解”的比喻：看懂一盘象棋——先找到规则，再解释具体走法。
  \item 规则对不对，有三种检验：简单系统可精确预测、由规则推出更一般的法则、用粗略近似把握意图。
  \item 已知规则最有趣之处，是它在“新场合”失效的地方——那是发现新规则的入口。
  \item 物理学史是一连串综合：热与力学、电与磁与光（电磁场）、化学与原子行为（量子化学）；本章最终追问：事物由什么构成、有多少基本元素。
\end{itemize}
\end{summarybox}


% ============================================================
% §2-3 量子物理学
% ============================================================
% 样张未收录 §2-2，单独编译时手动把节号补到 3，保持与全书一致的「2.3」
\setcounter{section}{2}
\section{量子物理学}

\runin{经典物理的两道裂缝} 说明了电磁场概念和电磁场能传送波后，我们很快就认识到，这些波的行为实际上十分奇怪，看起来完全不像波。在频率较高时，它们的行为更像\uwave{粒子}！正是在1920年后发展起来的\term{量子力学}解释了这种奇怪的行为。在1920年之前，爱因斯坦已改变了把空间看作是三维空间、把时间看成是单独存在的这种图像。他首先把它们组合在一起，并称之为\term{空－时}，然后又进一步用弯曲的空－时来描绘万有引力。这样，宇宙的“舞台”就变为空－时，而万有引力则大概是空－时的一种变态。

以后，人们又发现有关原子运动的规则也是有问题的：在原子世界中，“惯性”与“力”的力学法则是不正确的——牛顿定律已不再成立。人们反而发现小尺度范围内事物的行为与大尺度范围内事物的行为没有任何相似之处。这给物理学造成困难——但又十分有趣。之所以困难，是由于事物在小尺度范围内的表现如此“反常”，我们对之没有直接的经验。在这里事物的表现完全不像我们所知道的任何事情，因而除了用解析的方式，用任何其他方法都不可能描写这种习性。这的确是困难的，需要做大量的想象。

\runin{测不准原理} 量子力学中有许多看法。首先，一个粒子既有确定的位置也有确定的速度，这种概念已被抛弃，那是不正确的想法。表明经典物理是怎样不正确的一个例子是，在量子力学中有这样一条定则：不可能既知道某个粒子在什么地方，又知道它运动得多块。\keyidea{动量的不确定性与位置的不确定性是并协的，二者的乘积是常数。}我们可以把这条定律写成$ \Delta x\,\Delta p\geq\hbar/2 $，在以后将会更详尽的解释它。

这条定则解释了这样一个十分神秘的佯谬：即如果原子是由正负电荷所构成，那么为什么负电荷不是简单地位于正电荷的顶端（它们彼此是吸引的），从而彼此靠拢以至于完全抵消？\uwave{为什么原子如此庞大}？为什么原子核在中心，而其周围环绕着一些电子？

起先曾认为原子核很大，但事实并非如此，它是\uwave{非常小}的。一个原子的直径约为$ 10^{-8} $厘米，一个原子核的直径约为$ 10^{-13} $厘米。如果我们有一个原子，为了看到原子核，就要把整个原子放大到一个大房间大样大。这时原子核才刚刚是一个可以用眼睛分辨出来的斑点，但是几乎原子所有的重量都集中在这个无比小的原子核上。是什么理由使电子没有直接落入原子核呢？正是上述的原理。

如果电子在原子核里出现，我们就会精确地知道它们的位置，而\term{测不准原理}则要求它们具有很大的（不过是不确定的）动量，即很大的动能。电子具有这样大的能量就要脱离原子核，这些电子作出了让步：由于不确定性，它们为自己留下一个狭小的空间，于是以由这个定则所决定的最小的运动晃动着。

（记得我们曾经说过，当晶体冷却到绝对零度时，原子并没有停止运动，它们仍然在晃动，为什么？如果它们停止运动，我们就能知道它们在什么地方，而且它们不运动，这就违反了测不准原理：我们不能既知道它们在哪里，又知道它们以什么速度运动。所以它们必须在那里不断地摆动！）

\runin{不可精确预言} 另一个由量子力学带来的在科学的观念和哲学方面最有趣的变化是，在任何情形下要想精确地预言会发生什么事都是不可能的。比如我们有可能使一个原子处于准备发光的状态，在原子发光时，可以利用探测\term{光子}的方法进行测量（这一点我们马上就要讲的），但是，我们无法预计它将在什么时候发光，或者在有几个原子的情况下，究竟哪一个原子将发光。你们可能说，这是由于某种我们还没有足够仔细观察过的内部“转轮”在起作用。然而，这里根本没有什么内部的转轮，按照我们今天的理解，大自然的表现是这样的：\keyidea{根本不可能精确地预言在一定的实验中究竟会发生什么事情。}

这是一件糟透了的事；事实上，哲学家曾声称：科学所必需的基本东西之一就是，每当你安排了同样的条件时，那么发生的必定是同一件事。但是，这完全不正确，它并不是科学的基本条件。事实上是所发生的并不是同一件事，我们所能得到的只是发生一些什么的统计平均。不过，科学并没有完全崩溃。

\begin{optional}顺便地说，哲学家们讲了一大套科学之绝对必需是什么，但就像人们所看到的那样，这些总是相当天真的，甚至还是错误的。例如，某个哲学家宣称对科学的成就来说十分重要的是，如果同一个实验先在某处，比如说在斯德哥尔摩做；然后在另一处，比如说在基多（南美厄瓜多尔首都——译者注）做，那么必定会出现同样的结果。这纯粹是一派胡言。

对科学来说，这并不是必然的：它可能是一个经验事实，但不是必然的情况。比如有一个实验室在斯德哥尔摩观察天空，这时会看到北极光，如果在基多则看不到这种现象，这就是出现了不同的情况。“但是”，你会说：“这是一件与外部情况有关的事，如果你把自己关在斯德哥尔摩的一个房间里，拉下窗帘的话，那么会发生什么差别吗？”肯定会。假如我们在一个方向接头上挂一个摆，让它开始摆动，它就会差不多在一个平面里摆动，但也并不完全如此。在斯德哥尔摩，平面会缓慢地转动着，但是在基多就不会。在那里，窗帘也是垂下的。这件事的发生并没有引起科学的毁灭。\end{optional}

科学的基本假设，它的基本哲学观念是什么呢？

\pullquote{我们在第一章里讲到过：实验是任何观念的正确性的唯一试金石。}

假如结果是在基多所作的大多数实验与在斯德哥尔摩所作的实验效果一样，那么这“大多数实验”就可用来提出某种一般性的定律，至于对那些效果不同的实验，我们就将说：“这是由于斯德哥尔摩周围的环境不同所引起的”。我们将能想出一些办法来概括实验结果，而没有必要在事先就被告诫说，这些办法看起来象什么。假如有人告诉我们说，同样的实验总是产生同样的结果，这固然很好。但是当我们试了一下后，发现并非如此，因而结论的确就是并非如此。我们正是必须相信自己所看到的，然后才能借助于实际的经验来形成我们的一切其他观念。

\runin{波粒二象性} 现在让我们回到量子力学和基本物理上来。当然，我们在此刻还不能详细叙述量子力学的原理，因为它们是颇难理解的。我们将假定它们成立，然后叙述一下某些结果。其中一个是，我们通常视作为波的那些事物也具有粒子的习性，而粒子则具有波的习性。\keyidea{实际上，每一种事物的行为都是一样的，不存在波和粒子的区别。}

这样，量子力学就将场的概念以及场的波与粒子统一起来。的确，频率低时，现象的场的方面比较明显，或者说作为根据日常经验的近似描写时比较有用。但当频率增加时，现象的粒子方面对于我们通常用来作为测量用的仪器来说更为明显。实际上，虽然我们提到过许多频率，但目前还没有探测到任何直接涉及频率在每秒$10^{12}$周以上的现象，我们只是在假定了量子力学的\term{波粒二象性}概念是正确之后，根据有关规则从粒子的能量来推断出这些较高的频率的。

\runin{量子电动力学} 于是，我们对电磁相互作用有了新的见解。我们把一种新的粒子加入到电子、质子及中子的行列，这种新的粒子称为光子。新的电子质子相互作用的见解称为\term{量子电动力学}，它就是电磁理论，不过其中的一切在量子力学上都是正确的。\keyidea{这是光和物质，或电场与电荷之间的相互作用的基本理论，就物理学来说它是我们最伟大的成就。}

比如，从量子电动力学可以得出所有已知的电学、力学和化学定律：

\begin{itemize}
  \item 弹子碰撞的定律，
  \item 导线在磁场中运动的定律，
  \item 一氧化碳的比热，
  \item 霓虹灯的色彩，
  \item 盐的密度，
  \item 以及氢气与氧形成水的反应等，
\end{itemize}

全部都是这一理论的推论。

所有这些细节，如果简单到能使我们运用近似方法的话，都可以得出，这实际上当然不可能。不过，我们总能对发生的事多少有所理解。目前，在原子核外面还没有发现量子电动力学定律有什么例外，对于原子核我们不知道是否会有例外，因为对于核内的过程我们简直还不太清楚。

\runin{成就与边界} 这样，在原则上，量子电动力学是一切化学以及生命的理论——如果生命最后归结为化学，因而也就归结为物理的话（因为化学本身已经归结为物理，涉及化学中的那部分物理早就知道了！）。不仅如此，量子电动力学这个伟大的理论还预言了许多新的事实。首先，它说明了甚高能光子、γ 射线等等的性质。它还预言了另一个十分出乎意外的事：除电子外，还应当有同样质量、但带有正电荷的称为\term{正电子}的粒子，并且这两种粒子碰在一起时，会彼此湮没而放出光或 γ 射线（其实，光与 γ 射线完全是一回事，只是频率不同而已）。这件事情的推广——即对每个粒子总有一个\term{反粒子}——现在知道是正确的。

电子的反粒子有另一个名称，即正电子。但其他大多数反粒子，就称反某某子，如反质子、反中子。在量子电动力学中，提出了两个基本数据——电子质量与电荷，所有世界上其他的数被认为可以从这两个数据推导出来。实际上，这不完全正确，因为化学还有一整套数据，它告诉我们原子核是多重，这就把我们引导到下一部分内容中去了。

\begin{conceptmap}[示意：量子力学的三大观念]
\centering
\begin{tikzpicture}[font=\small,>=stealth,
  box/.style={draw,rounded corners=2pt,align=center,text width=3.15cm,
              inner sep=3pt,font=\footnotesize,fill=white},
  arr/.style={->,thick,gray!65}]
\node[box] (q) at (0,3.2) {\textbf{量子力学}\\(1920 年后)};
\node[box] (a) at (-4.3,1.2) {\textbf{1. 测不准原理}\\位置与动量不可同时确定};
\node[box] (b) at (0,1.2)   {\textbf{2. 不可精确预言}\\只能给出统计平均};
\node[box] (c) at (4.3,1.2) {\textbf{3. 波粒二象性}\\波与粒子没有区别};
\node[box,draw=blue!55] (d) at (0,-1.5) {\textbf{量子电动力学}\\统一光、电与物质};
\draw[arr] (q.south) -- (a.north);
\draw[arr] (q.south) -- (b.north);
\draw[arr] (q.south) -- (c.north);
\draw[arr] (a.south) to[out=-70,in=180] (d.west);
\draw[arr] (b.south) -- (d.north);
\draw[arr] (c.south) to[out=-110,in=0] (d.east);
\end{tikzpicture}

\figsource{按 §2-3 正文归纳绘制，非原书插图}
\end{conceptmap}

\begin{summarybox}[§2-3 小结]
\begin{itemize}
  \item 舞台变了：相对论把空间与时间合成空－时；小尺度上牛顿定律失效，日常直觉不再可靠。
  \item 测不准原理：位置与动量不能同时确定；它解释了原子为何“庞大”、电子为何不落入原子核。
  \item 量子力学放弃“精确预言”，只给统计平均；但科学没有崩溃——检验观念的标准仍然是实验。
  \item 波粒二象性：波与粒子不是两种东西；量子力学把场的波与粒子统一起来。
  \item 量子电动力学的成功与边界：它统一了光、电与物质，但电子质量与电荷仍是输入数据，核内规律尚不清楚。
\end{itemize}
\end{summarybox}
